\chapter{Structural Results}\label{chap:structural-results}

This chapter shall be devoted to looking at some structural results on arithmetic circuits. 
This would help us understand the relevance of shallow circuits in the context of proving lower bounds for arithmetic circuits of arbitrary depth. 

\section{Log-product representations for formulas}

The following structural lemma shows that any multilinear formula can be converted in to a small sum of \emph{log-product} polynomials. 
The techniques of the following lemma can also be used in other settings with minor modifications, and we shall encounter a different version of this lemma later as well.
These normal forms is from the survey of Shpilka and Yehudayoff~\cite{sy}, and also from the result of \Hrubes and Yehudayoff~\cite{HY11a}. 

\begin{definition}\label{defn:mult-logproduct}
  A multilinear polynomial $f\in \F[X]$ is called a \emph{multilinear log-product} polynomial if $f = g_1\dots g_t$ and there exists a partition of variables $X = X_1 \sqcup \dots \sqcup X_t$ such that
  \begin{itemize}
  \item $g_i \in \F[X_i]$ for all $i \in [t]$.
  \item $\pfrac{1}{3}^i |X| \leq |X_i| \leq \pfrac{2}{3}^i |X|$ for all
    $i$, and $|X_t| = 1$.
  \end{itemize}
\end{definition}

\begin{lemma}\label{lem:mult-logproduct}
  Let $\Phi$ be a multilinear formula of size $s$ computing a polynomial $p$. 
Then $p$ can be written as a sum of $(s+1)$ log-product multivariate polynomials.
\end{lemma}
\begin{proof}
  Similar to \autoref{lem:formula-depth-reduction}, let $v$ be a node in $\Phi$ such that set of variables $X_v$ that it depends on satisfies $\frac{|X|}{3} \leq \abs{X_v} \leq \frac{2|X|}{3}$. 
If $\Phi_v$ is the polynomial computed at this node, then $f$ can be written as
  $$
  f \spaced{=} \Phi_v \cdot g_1 + \Phi_{v=0} \quad\text{for some $g_1 \in \F[X\setminus X_v]$}.
  $$
  where $\Phi_{v=0}$ is the formula obtained by replacing the node $v$ by zero. 
Note that the subtree at the node $v$ is completely disjoint from $\Phi_{v=0}$. 
Hence the sum of the sizes of $\Phi_v$ and $\Phi_{v=0}$ is at most $s$. 
Hence, $g_1 \in \F[X\setminus X_v]$ and $\frac{|X|}{3} \leq \abs{X \setminus X_v} \leq \frac{2|X|}{3}$. 
Inducting on the formulas $\Phi_v$ and $\Phi_{v=0}$ gives the lemma.
\end{proof}

There are many variants of this formula depending on how you pick the node $v$ in the proof above.
Here is another variant that find a node $v$ based on \emph{degree} rather than the number of variables that it depends on.
We state it here without proof but it follows exactly as in the lines of \autoref{lem:mult-logproduct}.

\begin{lemma}\label{lem:hom-logproduct}
  Let $\Phi$ be a homogeneous formula of size $s$ computing a polynomial $p$ of degree $d$.
Then $p$ can be written as a sum of $(s+1)$ \emph{log-product} polynomials, that is,
\[
p \spaced{=} f_1 + \cdots + f_{r}\quad\quad \text{with $r \leq s+1$}
\]
where for each $i \in [r]$, we have $f_i = f_{i1} \cdots f_{i\ell}$ satisfying
\begin{itemize}
\item each $f_{ij}$ is homogeneous,
\item $(1 / 3)^j \cdot d \leq \deg(f_{ij}) \leq (2/ 3)^j \cdot d$,
\item $f_{i\ell} = 1$.
\end{itemize}
In particular, each $f_i$ factors into $\Omega(\log d)$ non-trivial factors of geometrically decreasing degrees. 

Furthermore, if $\Phi$ was a multilinear formula to begin with, then so is the expression on the RHS. 
\end{lemma}

Yet another variant of this trick is the following.


%%% Local Variables: 
%%% mode: latex
%%% TeX-master: "fancymain"
%%% End: 
